\scanpage{91}

\deleted{(\unsure{indp. de $i$})}

La donnée d'un tel $c_\Omega$ équivaut à une restriction du Groupe structural de $\underline{\mathrm{Conn}}_i$ \deleted{\unsure{$\mathscr L$|$\varepsilon$}} à $\deleted{\unsure{\omega}}\Omega\otimes\underline\omega$, ou encore (\unsure{indépendamment|indépendant} de $i$) la donnée d'un $c_\Omega$ pour lequel ($\ast$) est lin et ($\ast\ast$) est \unsure{vér.|vérifiée|vraie}, équivaut à une restriction du Groupe structural $\unsure{\mathfrak g|\mathcal G}\otimes\omega$ de $\underline{\mathrm{Conn}}$ à $\Omega\otimes\underline\omega$..

On s'intéresse au cas où \deleted{\unsure{$c_i$|$\underline{\mathrm{Conn}}_i$}} $c_\Omega$ est tel que la composée
\[ (-2i):\Omega\longrightarrow\mathscr L\otimes\underline\omega\unsure{\relbar|\longrightarrow}\underline\omega \]
soit un \uline{iso}. \inserted{(\uline{ce qui implique 2 inv.})} Alors \deleted{\unsure{le}} le triplet

$(\mathfrak g,\mathscr L,c_\Omega)$ (\deleted{ou} ce qui revient au $=$, $(\mathscr L\supset\Omega,c_\Omega)$) est défini à isom unique près par $X/S$ (et existe ssi $\underline\omega=\underline\Omega^1_{X/S}$ est inversible) comme
\[ \mathscr L\simeq\mathcal P^1_{X/S}(\underline\omega)\otimes\underline\omega^{-1}\quad(\text{donc }\mathscr L\otimes\underline\omega=\mathcal P^1_{X/S}(\underline\omega)) \]

et $c_\Omega:\Omega\simeq\underline\omega\longrightarrow\mathscr L\otimes\underline\omega\simeq\mathcal P^1_{X/S}(\underline\omega)$ n'est autre que $d^1_{X/S,\deleted{\unsure{\omega}}\underline\omega}$ --

% archive sheet number: 90
